\documentclass[12pt,a4paper]{article}

\usepackage[italian]{babel}
\usepackage[utf8]{inputenc}
\usepackage[T1]{fontenc}
\usepackage{geometry}
\usepackage{graphicx}
\usepackage{setspace}
\usepackage{titlesec}

\geometry{margin=2.5cm}
\setstretch{1.2}

\titleformat{\section}{\large\bfseries}{\thesection.}{0.5em}{}
\titleformat{\subsection}{\normalsize\bfseries}{\thesubsection.}{0.5em}{}

\title{\textbf{Programma Elettorale}\\
\vspace{1cm}

\includegraphics[width=0.2\textwidth]{assets/logo_redbg.png}\\[1em]
\vspace{0.5em}
\large \textbf{Componente Connessa}}
\author{Eduardo Perreca, Riccardo Pasquino, Adriana Carollo, \\ Lorenzo Medde, Gabriele Fioco, Elijah Lisa Bassetti}
\date{}

\begin{document}

\maketitle

\section*{Introduzione}

Componente Connessa nasce dall’esigenza \textbf{concreta} di unire le diverse realtà di rappresentanza presenti all’interno del Dipartimento di Informatica in un unico progetto condiviso. Ci proponiamo come una lista \textbf{apartitica} e \textbf{indipendente}, focalizzata \textbf{esclusivamente} sull’ascolto e sulla risoluzione dei problemi reali riscontrati dagli studenti nel Dipartimento di Informatica e nella Facoltà di Scienze e Tecnologie.
\\ Il nostro gruppo è composto da \textbf{studenti motivati}, \textbf{rappresentanti con esperienza}, pronti a impegnarsi attivamente e con attenzione verso temi fondamentali quali la \textbf{qualità della didattica}, l’\textbf{organizzazione degli spazi}, il \textbf{benessere studentesco} e l’\textbf{accessibilità dei servizi}.
\\ Ci presentiamo come unica lista di \textbf{rappresentanza dipartimentale} del Dipartimento di Informatica, con l’obiettivo di garantire una \textbf{voce unitaria}, forte e coerente. Allo stesso tempo, estendiamo la nostra candidatura anche nella \textbf{facoltá} di \textbf{Scienze e Tecnologie}, per \textbf{rappresentare} e \textbf{tutelare} gli interessi dei dipartimenti che rimangono nell’area di Città Studi, a seguito del trasferimento a \textit{MIND}, promuovendo una collaborazione più \textbf{efficace} e \textbf{puntuale}.

\newpage
\tableofcontents

\section{Metodo}

Componente Connessa si impegna a operare attraverso:
\begin{itemize}
    \item un approccio inclusivo, apartitico e indipendente, aperto a chiunque voglia contribuire attivamente alla rappresentanza, tramite segnalazioni e incontri;
    \item l'ascolto continuo e proattivo degli studenti, attraverso momenti di confronto e spazi dedicati al dialogo;
    \item la raccolta e valutazione costante dei feedback, con particolare attenzione ai singoli corsi e alle loro criticità;
    \item la divulgazione di informazioni che permettano agli studenti di acquisire consapevolezza e autonomia riguardo il proprio percorso di studi e ciò che lo circonda, con particolare attenzione al rendere le informazioni accessibili sia agli studenti italiani che internazionali.
\end{itemize}

\section{Le nostre proposte}
\subsection{Questionari Post-Esame}
Nell’ultimo anno abbiamo promosso l’introduzione dei questionari post-esame per affiancare gli strumenti di valutazione esistenti e ottenere un’analisi più completa dei corsi. Proponiamo di consolidarne e ampliarne l’utilizzo, lavorando per un loro riconoscimento formale e per l’integrazione nei processi di valutazione della didattica, anche tramite la Commissione Paritetica, al fine di raccogliere dati più rappresentativi e migliorare concretamente la qualità dei corsi.\\Crediamo che gli insegnamenti non debbano essere valutati solo tramite le lezioni, ma anche considerando la correttezza e la coerenza dei metodi di valutazione adottati, migliorando così la qualità complessiva della didattica.

\subsection{Trimestri Magistrale}
La recente introduzione dei nuovi modelli organizzativi della didattica in magistrale ha portato a una maggiore concentrazione dei corsi, generando difficoltà nella gestione del carico di studio, sovrapposizioni e una fruizione meno efficace delle lezioni.
\\Proponiamo di intervenire sulle modalità di valutazione, promuovendo un sistema più distribuito nel tempo che non si basi esclusivamente sull’esame finale. In particolare, sosteniamo l’introduzione di prove intermedie, progetti, lavori di gruppo e altre forme di valutazione continua, così da rendere il percorso più sostenibile e permettere agli studenti di seguire più corsi con maggiore equilibrio.
\\Parallelamente, riteniamo fondamentale raccogliere dati concreti sull’impatto di questo modello, attraverso sondaggi e strumenti di monitoraggio, per comprendere il numero di esami sostenuti e le principali criticità. Queste informazioni saranno essenziali per portare in sede decisionale una valutazione informata sull’efficacia del sistema e proporre eventuali miglioramenti o modifiche.
\subsection{Percorso Videogames}
Il percorso magistrale in ambito videogiochi presenta attualmente diverse criticità, in particolare legate al carico di lavoro richiesto rispetto ai crediti formativi assegnati. Molti corsi risultano infatti eccessivamente impegnativi rispetto ai 6 CFU previsti, creando difficoltà nel bilanciamento del percorso di studi.
\\Proponiamo di avviare un confronto con il corpo docente per rivedere l’organizzazione dei corsi e delle modalità di valutazione, con l’obiettivo di renderle più coerenti con il carico formativo previsto. Riteniamo fondamentale garantire una maggiore proporzionalità tra impegno richiesto e crediti, migliorando così la sostenibilità e l’efficacia del percorso.

\subsection{Aule di Celoria e Venezian}
Le aule situate in via Celoria e via Venezian presentano attualmente diverse criticità, in particolare per quanto riguarda le infrastrutture e i servizi disponibili agli studenti. Tra i problemi più rilevanti si evidenzia la carenza di prese elettriche adeguate, che rende difficile seguire le lezioni in modo efficace, soprattutto in un contesto come quello informatico.
\\Alla luce anche del progressivo trasferimento verso il polo MIND, emerge una limitata attenzione istituzionale verso il miglioramento di questi spazi. Riteniamo tuttavia che sia possibile intervenire concretamente già a livello dipartimentale.
\\Proponiamo quindi di avviare azioni mirate per migliorare la vivibilità delle aule, a partire da interventi semplici ma impattanti come l’installazione di ulteriori prese elettriche e il potenziamento delle dotazioni esistenti, con l’obiettivo di garantire un ambiente di studio più adeguato alle esigenze degli studenti.\\Proponiamo inoltre di rendere libero l'accesso alle aule in cui non si stanno svolgendo lezioni, aumentando così gli spazi disponibili e accessibili allo studio.

\subsection{Registrazioni delle Lezioni}
Riteniamo che la disponibilità di registrazioni delle lezioni rappresenti uno strumento fondamentale per migliorare la fruibilità degli insegnamenti. In particolare, esse risultano utili nei casi di sovrapposizioni tra corsi o difficoltà organizzative, permettendo agli studenti di accedere ai contenuti in maniera più flessibile.
\\Siamo convinti che l’introduzione delle registrazioni non riduca la partecipazione in presenza, ma costituisca invece un supporto complementare allo studio.
\\Proponiamo quindi di avviare una sperimentazione, a partire dai corsi opzionali sia triennali che magistrali, con l’obiettivo di estendere progressivamente la disponibilità di materiale didattico accessibile in modalità asincrona e migliorare complessivamente l’esperienza di apprendimento.

\subsection{Tutoraggio}

Riteniamo fondamentale potenziare le attività di tutoraggio, in particolare per gli insegnamenti caratterizzati da un'elevata difficoltà e da una bassa percentuale di superamento dell'esame.
\\Proponiamo un incremento del numero di tutoraggi disponibili, valorizzando il ruolo degli studenti tutor, al fine di offrire un supporto più accessibile e mirato.

\subsection{Giornata dei Tirocini}

Riteniamo fondamentale supportare gli studenti in una scelta più consapevole del tirocinio e del percorso di tesi, sia triennale che magistrale, migliorando il dialogo con i docenti e i gruppi di ricerca.
\\Proponiamo l’organizzazione di una o più giornate dedicate, durante le quali i laboratori possano presentare i propri progetti, le aree di ricerca e le opportunità disponibili. Questo permetterebbe agli studenti di orientarsi meglio tra le diverse possibilità, riducendo l’incertezza nella scelta e aumentando la consapevolezza del proprio percorso.
\\Un’iniziativa di questo tipo contribuirebbe inoltre a rafforzare il collegamento tra studenti e attività di ricerca, con potenziali benefici anche per chi intende proseguire verso percorsi più avanzati come il dottorato.

\subsubsection{Gruppi Studenteschi}
Riteniamo utile per tutti gli studenti avere accesso alle iniziative organizzate da parte dei gruppi studenteschi di area informatica, favorendone le attività e comunicazioni.
In particolare, proponiamo l'inserimento sul sito del dipartimento dei collegamenti alle loro pagine web e il tempestivo aggiornamento degli eventi con le proposte dei gruppi, anche quando non ancora completamente finalizzate.
\\Crediamo che l'università non sia solo un luogo di studio individuale: siamo convinti che una maggiore partecipazione alle attività proposte dai gruppi studenteschi possa creare un ambiente ideale per il confronto, la collaborazione e la condivisione di idee.

\subsection{Internazionalizzazione}
Riteniamo fondamentale aumentare la consapevolezza delle opportunità di internazionalizzazione offerte dall'Ateneo. Molti studenti non sono pienamente informati sui programmi disponibili, come Erasmus e altri scambi internazionali, o sulle modalità per accedervi.
\\Proponiamo di migliorare la comunicazione e la visibilità di queste opportunità, attraverso una maggiore presenza di informazioni sui siti didattici di riferimento, l'organizzazione di incontri informativi dedicati e un supporto più diretto da parte dei rappresentanti, al fine di accompagnare gli studenti in un percorso universitario più aperto e internazionale.


\section*{Conclusione}
Il nostro obiettivo è costruire una comunità più connessa, equa e innovativa. Con il contributo di tutti vogliamo trasformare le idee in azioni concrete.
\\Per questo, il 22 e 23 marzo vota Componente Connessa alle elezioni in Dipartimento di Informatica e di Comitato di direzione della Facoltá di Scienze e Tecnologie.
\vfill


\vspace{1em}

\begin{center}
    \includegraphics[width=0.2\textwidth]{assets/logo_redbg.png}
\end{center}

\vspace{1em}

\begin{center}
\textit{Componente Connessa --- connettere per rappresentare, rappresentare per unire.}
\end{center}
\end{document}